% Part: many-valued-logic
% Chapter: three-valued-logics
% Section: kleene

\documentclass[../../../include/open-logic-section]{subfiles}

\begin{document}

\olfileid{mvl}{thr}{skl}

\olsection{Kleene 逻辑}

Stephen Kleene（斯蒂芬·克林）引入了两种三值逻辑，其动机源于这样一种逻辑：其中真值被视为计算过程的结果。一个过程可能输出 $\True$ 或 $\False$，但也可能无法终止。在后一种情况下，对应的真值是未定义的，用真值~$\Undef$ 表示。

要计算命题~$!A$ 的否定，首先计算 $!A$ 的值，然后返回其相反的值。如果 $!A$ 的计算不终止，那么整个计算过程也不会终止：因此 $\Undef$ 的否定仍是 $\Undef$。

要计算合取式 $!A \land !B$，有两种做法：可以先计算~$!A$，再计算~$!B$，若两者结果均为~$\True$ 则合取式为 $\True$，否则为 $\False$。如果其中任何一个计算未能停机，整个过程亦然。因此在这种情况下，只要有一个合取支未定义，合取式也就未定义。析取式的情况与此类似。

然而，如果我们能并行地求值 $!A$ 和 $!B$，结果会更好。此时，如果两个过程中有一个停机并返回 $\False$，我们就可以停止计算，因为该合取式的值必定为假。所以在这种情况下，只要有一个合取支为假，即使另一个合取支未定义，合取式也为假。类似地，当并行计算一个析取式时，一旦两个析取支中有一个的计算返回真，我们就可以停止：此时析取式必定为真。因此在这种情况下，即使其中一个支命题未定义，我们也能够知道复合命题的结果。在这种解释下，我们可以把 $\Undef$ 理解为“未知”而非“未定义”。

这两种解释分别导出了 Kleene 的强逻辑与弱逻辑。条件句被定义为与 $\lnot !A \lor !B$ 等价。

\begin{defn}
\emph{强 Kleene 逻辑}~$\LogKs$ 由以下矩阵定义：
\begin{enumerate}
  \item 标准的命题语言 $\Lang L_0$，包含
  $\lnot$、$\land$、$\lor$、$\lif$。
  \item 真值集合为 $V = \{\True, \Undef, \False\}$。
  \item $\True$ 是唯一的指定值，即 $V^+ = \{\True\}$。
  \item 真值函数由以下表格给出：
  \begin{center}
    \begin{tabular}{c|c}
      $\tf{\lnot}$ & \\
      \hline
      $\True$ & $\False$ \\
      $\Undef$ & $\Undef$ \\
      $\False$ & $\True$
    \end{tabular}
    \quad
    \begin{tabular}{c|ccc}
      $\tf{\land}[\LogKs]$ & $\True$ & $\Undef$ & $\False$ \\
      \hline
      $\True$ & $\True$ & $\Undef$ & $\False$ \\
      $\Undef$ & $\Undef$ & $\Undef$ & $\False$\\
      $\False$ & $\False$ & $\False$ & $\False$
    \end{tabular}
    \\[2ex]
    \begin{tabular}{c|ccc}
      $\tf{\lor}[\LogKs]$ & $\True$ & $\Undef$ & $\False$ \\
      \hline
      $\True$ & $\True$ & $\True$ & $\True$ \\
      $\Undef$ & $\True$ & $\Undef$ & $\Undef$ \\
      $\False$ & $\True$ & $\Undef$ & $\False$
    \end{tabular}
    \quad
    \begin{tabular}{c|ccc}
      $\tf{\lif}[\LogKs]$ & $\True$ & $\Undef$ & $\False$ \\
      \hline
      $\True$ & $\True$ & $\Undef$ & $\False$ \\
      $\Undef$ & $\True$ & $\Undef$ & $\Undef$  \\
      $\False$ & $\True$ & $\True$ & $\True$
    \end{tabular}
  \end{center}
\end{enumerate}
\end{defn}

\begin{defn}
\emph{弱 Kleene 逻辑}~$\LogKw$ 由以下矩阵定义：
\begin{enumerate}
  \item 标准的命题语言 $\Lang L_0$，包含
  $\lnot$、$\land$、$\lor$、$\lif$。
  \item 真值集合为 $V = \{\True, \Undef, \False\}$。
  \item $\True$ 是唯一的指定值，即 $V^+ = \{\True\}$。
  \item 真值函数由以下表格给出：
  \begin{center}
    \begin{tabular}{c|c}
      $\tf{\lnot}$ & \\
      \hline
      $\True$ & $\False$ \\
      $\Undef$ & $\Undef$ \\
      $\False$ & $\True$
    \end{tabular}
    \quad
    \begin{tabular}{c|ccc}
      $\tf{\land}[\LogKw]$ & $\True$ & $\Undef$ & $\False$ \\
      \hline
      $\True$ & $\True$ & $\Undef$ & $\False$ \\
      $\Undef$ & $\Undef$ & $\Undef$ & $\Undef$\\
      $\False$ & $\False$ & $\Undef$ & $\False$
    \end{tabular}
    \\[2ex]
    \begin{tabular}{c|ccc}
      $\tf{\lor}[\LogKw]$ & $\True$ & $\Undef$ & $\False$ \\
      \hline
      $\True$ & $\True$ & $\Undef$ & $\True$ \\
      $\Undef$ & $\Undef$ & $\Undef$ & $\Undef$ \\
      $\False$ & $\True$ & $\Undef$ & $\False$
    \end{tabular}
    \quad
    \begin{tabular}{c|ccc}
      $\tf{\lif}[\LogKw]$ & $\True$ & $\Undef$ & $\False$ \\
      \hline
      $\True$ & $\True$ & $\Undef$ & $\False$ \\
      $\Undef$ & $\Undef$ & $\Undef$ & $\Undef$  \\
      $\False$ & $\True$ & $\Undef$ & $\True$
    \end{tabular}
  \end{center}
\end{enumerate}
\end{defn}

\begin{prop}
  $\LogKs$ 和 $\LogKw$ 没有重言式。
\end{prop}

\begin{proof}
  如果对所有命题变元~$p$，都有 $\pAssign v(p) = \Undef$，
  那么任何公式~$!A$ 的真值将是 $\pValue v(!A) =
  \Undef$，因为在两种逻辑中 \[
    \tf{\lnot}(\Undef) = \tf{\lor}(\Undef, \Undef) = \tf{\land}(\Undef,
  \Undef) = \tf{\lif}(\Undef, \Undef) = \Undef
  \] 均成立。由于在 $\LogKs$ 或
  $\LogKw$ 中 $\Undef \notin V^+$，
  因此在该赋值下，$!A$ 的值不是指定值。
\end{proof}

虽然弱和强 Kleene 逻辑都没有重言式，但它们具有非平凡的后承关系。

\begin{prob}
  下列后承关系在（a）强和（b）弱 Kleene 逻辑中哪些成立？对每个关系给出真值表。
  \begin{enumerate}
    \item $p, p \lif q \Entails q$
    \item $p \lor q, \lnot p \Entails q$
    \item $p \land q \Entails p$
    \item $p \Entails p \land p$
    \item $p \Entails p \lor q$
  \end{enumerate}
\end{prob}

Dmitry Bochvar（德米特里·博赫瓦尔）将 $\Undef$ 解释为“无意义”，并试图通过规定悖论性语句（如说谎者悖论）取真值~$\Undef$ 来解决此类悖论。他引入了一种逻辑，其本质上是弱 Kleene 逻辑增加了额外联结词的扩展，其中两个是“外部否定”和“是未定义的”算子：
\begin{center}
  \begin{tabular}{c|c} 
    $\tf{\sim}$ & \\ 
    \hline  
    $\True$ & $\False$ \\ 
    $\Undef$ & $\True$ \\
    $\False$ & $\True$ 
  \end{tabular}
\quad
  \begin{tabular}{c|c} 
    $\tf{+}$ & \\ 
    \hline  
    $\True$ & $\False$ \\ 
    $\Undef$ & $\True$ \\
    $\False$ & $\False$ 
  \end{tabular}
\end{center}

\begin{prob}
  你能否用 $\lnot$ 和~$+$ 在 Bochvar 逻辑中定义 $\sim$？
  即，找出一个只含命题变元~$p$ 且不涉及 $\sim$ 的公式，使其总是取与~$\mathord{\sim}p$ 相同的真值。
  用真值表证明你是正确的。
\end{prob}


\end{document}